% Generated draft: requires mathematical and source review.
\documentclass[11pt]{article}
\usepackage[margin=25mm]{geometry}
\usepackage{fontspec}
\setmainfont{Latin Modern Roman}
\usepackage{amsmath,amssymb,mathtools,unicode-math}
\setmathfont{Latin Modern Math}
\usepackage{enumitem,needspace,xurl,hyperref}
\setlength{\parindent}{0pt}
\setlength{\parskip}{5pt}
\setlength{\emergencystretch}{3em}
\newcommand{\N}{\mathbb{N}}
\newcommand{\Z}{\mathbb{Z}}
\newcommand{\Q}{\mathbb{Q}}
\newcommand{\R}{\mathbb{R}}
\newcommand{\C}{\mathbb{C}}
\newcommand{\F}{\mathbb{F}}
\newcommand{\A}{\mathbb{A}}
\newcommand{\PP}{\mathbb{P}}
\newcommand{\E}{\mathbb{E}}
\newcommand{\eps}{\varepsilon}
\newcommand{\dd}{\,\mathrm{d}}
\begin{document}
\section*{UnsolvedMath ID: 30003542}
\textbf{Source problem code:} catalogue entry
\subsection{Definitions and mathematical statement}
The supplied record contains two separate questions. Its missing hypotheses prevent a single fully specified formulation; the following definitions make the targets explicit while identifying the unresolved choices.

\textbf{Closed curves of prescribed geodesic curvature.} To define ordinary signed geodesic curvature, let \((S,g)\) be a smooth oriented Riemannian surface diffeomorphic to \(S^2\). Let \(\nabla\) be its Levi--Civita connection, and let \(J\) denote rotation through \(\pi/2\) in each oriented tangent plane. For \(L>0\), a smooth unit-speed embedding
\[
\gamma:\R/L\Z\longrightarrow S
\]
has signed geodesic curvature
\[
k_g(s)=g\bigl(\nabla_{\dot\gamma(s)}\dot\gamma(s),J\dot\gamma(s)\bigr).
\]
Under the interpretation that the question concerns ordinary signed geodesic curvature on arbitrary smooth Riemannian two-spheres, its candidate formulation is:
\[
\text{For every }(S,g)\text{ as above and every }c\in\R,
\quad \exists L>0\ \exists\gamma:\R/L\Z\hookrightarrow S
\quad\text{such that}\quad
|\dot\gamma|_g=1,\qquad k_g(s)=c\ \text{for all }s.
\]
This is only a conditional interpretation. The source says that the sphere is ``in a magnetic field,'' but supplies neither the magnetic field as mathematical data nor its relation to the prescribed curvature. It also does not specify a metric, regularity, a signed or unsigned curvature convention, or the range of prescribed constants. Consequently, the displayed interpretation cannot be asserted to be the intended magnetic problem. The requirement that the curve be closed and embedded is explicit and must be retained.

\textbf{Multiplicity of min--max minimal hypersurfaces.} A minimal hypersurface in a Riemannian manifold is a hypersurface whose mean curvature vector vanishes. Whenever a min--max output can be represented as an integral varifold of the form
\[
V=\sum_{i=1}^{N}m_i|\Sigma_i|,
\qquad N\geq 1,\qquad m_i\in\Z_{>0},
\]
where the \(\Sigma_i\) are distinct connected smooth embedded minimal hypersurfaces and \(|\Sigma_i|\) denotes integration over \(\Sigma_i\) with unit density, the output has \emph{multiplicity one} if
\[
m_i=1\qquad\text{for every }i\in\{1,\ldots,N\}.
\]
The second supplied question asks whether the intended min--max construction always produces a minimal hypersurface with multiplicity one. In the representation above, this asks whether all component weights must equal one. The record does not specify the ambient manifolds or their dimensions, metric hypotheses, boundary conditions, sweepout or coefficient class, or the min--max construction. Nor does it specify whether the assertion concerns every possible output or the existence of a multiplicity-one choice of output for each admissible input. Smoothness and the displayed decomposition are conditions for this explanation, not hypotheses established by the source for every intended case. These missing choices prevent a fully quantified universal statement.

The record labels the combined problem \emph{partially solved}. Its supplied dated assessment states that Zhou and Zhu proved existence of a nontrivial closed constant-mean-curvature hypersurface for every prescribed constant in a closed ambient manifold, that nonzero min--max solutions have multiplicity one, and that the minimal multiplicity-one problem is known only under additional generic hypotheses. These are supplied literature claims, not independently verified conclusions here. Their precise hypotheses, the meaning of ``nonzero,'' and their implications for the embedded magnetic-curve question are not specified.
\subsection{Short English statement}
The record asks two separate questions: whether a two-sphere in a magnetic field has a closed embedded curve of every prescribed constant geodesic curvature, and whether minimal-hypersurface min--max constructions always yield multiplicity one. The magnetic data, curvature convention, and scope of the min--max assertion are unspecified. The supplied \emph{partially solved} label and literature assessment require independent review.
\subsection{Sources}
\begin{itemize}
\item Local UnsolvedMath registry: supplied record with permanent ID \(30003542\), exact source problem code \texttt{catalogue entry}, title ``Closed Magnetic Geodesics and Multiplicity-One Minimal Hypersurfaces,'' in the Oberwolfach Reports Open Problems collection. The project registry locations are \texttt{data/unsolvedmath/problems.json} and the supplementary \texttt{catalog/problems.json}; record membership in either file has not been independently checked here. The record supplies proposed year \(2017\) and no individual proposer attribution.
\item Supplied report citation: \emph{Partial Differential Equations (2018)}, Oberwolfach Reports, DOI \texttt{10.4171/owr/2017/35}, pp.\ 2165--2222. \url{https://doi.org/10.4171/owr/2017/35}. The report title year and DOI year are retained exactly as supplied.
\item Supplied literature reference: Zhou--Zhu, ``CMC min--max theory.'' \url{https://arxiv.org/abs/1707.08012}.
\item Supplied literature reference: Zhou, ``Multiplicity one for bumpy metrics.'' \url{https://annals.math.princeton.edu/2020/192-3/p03}.
\item Literature assessment date: \texttt{2026-08-21}. This is supplied dated curation, not an independently verified literature assessment. The curation labels the problem partially solved and reports the constant-mean-curvature existence and multiplicity claims described above; it explicitly requests independent verification of time-sensitive status and cited claims.
\item The registry assigns statement status \texttt{corrected\_verified} and says the clean statement was checked against the official Oberwolfach report. This is a supplied verification claim; no source visit or independent verification is asserted in this draft.
\end{itemize}
\end{document}
